\section{Experimental Setup}
\label{sec:setup}

\subsection{Dataset}
\label{sec:setup:dataset}

Beyond the headline counts already noted in Section~\ref{sec:intro}, FAR-Trans~\cite{sanzcruzado:fartrans} additionally provides 159{,}136 Sell transactions, with the 806 assets split across stocks, bonds, and mutual funds, and daily closing prices spanning January 2018 to November 2022. The MiFID II \texttt{riskLevel} field used throughout this paper takes values in \{Conservative, Income, Balanced, Aggressive\}, with \texttt{Predicted\_*} variants flagging regression-imputed bands. In addition, customer profiles also carry \texttt{customerType} and \texttt{investmentCapacity}, which enter the RQ2 controls in Section~\ref{sec:setup:stats}.

\subsection{Asset Risk Band Assignment}
\label{sec:setup:bandassignment}

Each asset is assigned to one of the four MiFID risk bands by a three-step hierarchical rule, applied in order.

\subsubsection{Step 1: subcategory metadata} Mutual funds are labelled directly from their subcategory: Money Market funds are treated as Conservative, Bonds as Income, Balanced as Balanced, and Equity or Large Cap as Aggressive. Bonds are labelled by issuer type: Government as Conservative, Corporate as Income, and any remaining bond subcategory as Income by default.

\subsubsection{Step 2: volatility quartiles for stocks} Stocks carry no subcategory tag we can use, so they fall through to a volatility rule. For each ISIN we compute the annualised standard deviation of daily log returns on a trailing 252-day window, then bucket all stocks into the four bands by the stock-only quartile thresholds $(q_1, q_2, q_3)$.

\subsubsection{Step 3: residual default} Any asset that falls through both prior steps is assigned to Balanced.

The resulting asset-band distribution across (Conservative, Income, Balanced, Aggressive) is $(190, 333, 105, 178)$ over $|A| = 806$ assets, which yields a band-conditional random baseline of $\pi(b) = (0.65, 0.78, 0.76, 0.35)$ under the formula defined in Section~\ref{sec:method:pck}.

\subsection{Temporal Split Protocol}
\label{sec:setup:splits}

Evaluation uses 69 monthly temporal splits on a trading-day grid between August 2019 and May 2022. For each split, training is all Buys before \texttt{time\_point}, and the test window is $(\texttt{time\_point}, \texttt{test\_end}]$, deduplicated against training and filtered to the training-asset universe. Eligible customers have at least one interaction in both training and test, while eligible assets have prices on both endpoints. In RQ4 (Section~\ref{sec:results:rq4}), per-split metric deltas are paired against the global LightGCN baseline to quantify the effect of each intervention component.

\subsection{Baselines and Hyperparameter Grids}
\label{sec:setup:baselines}

Table~\ref{tab:baseline-grids} reports the search grids and fixed values for the two FAR-Trans baselines. Both grids are swept with Ray Tune, with LightGCN tuned on average nDCG@10 and Random Forest on average ROI@10 across the 69 splits. Best-trial hyperparameters are available in the project repository (Section~\ref{sec:setup:reproducibility}).

\begin{table}[t]
\centering
\small
\begin{tabular}{lll}
\toprule
Parameter & Value & Role \\
\midrule
\multicolumn{3}{l}{\textbf{LightGCN} (tuned on average nDCG@10)} \\
\texttt{embedding\_dimension}   & $\{64, 128\}$         & grid \\
\texttt{number\_of\_layers}     & $\{2, 3\}$            & grid \\
\texttt{learning\_rate}         & $\{10^{-2}, 10^{-3}\}$ & grid \\
\texttt{weight\_decay}          & $10^{-5}$             & fixed \\
\texttt{keep\_probability}      & $0.6$                 & fixed \\
\texttt{number\_of\_epochs}     & $50$                  & fixed \\
\texttt{batch\_size}            & $1024$                & fixed \\
\midrule
\multicolumn{3}{l}{\textbf{Random Forest} (tuned on average ROI@10)} \\
\texttt{number\_of\_estimators} & $\{20, 30, 40, 50\}$  & grid \\
\texttt{max\_depth}             & $\{15, 25, 50\}$      & grid \\
horizon                         & 6 months              & fixed \\
\texttt{random\_state}          & $42$                  & fixed \\
\bottomrule
\end{tabular}
\caption{Baseline hyperparameter grids. ``grid'' rows are searched, ``fixed'' rows are held constant.}
\label{tab:baseline-grids}
\end{table}

The stratified PC-LightGCN inherits the best LightGCN trial's backbone hyperparameters. Two trials are evaluated, $\lambda = 0$ (BPR only, isolating the stratified architecture) and $\lambda = 1$ (coherence margin activated), on the same 69 splits.

\subsection{Statistical Inference}
\label{sec:setup:stats}

Three of the four research questions from Section~\ref{sec:intro} require statistical analysis beyond the headline metrics. RQ1 is descriptive and needs no inference test. The remaining three map to the following procedures.

\begin{itemize}
    \item \textbf{RQ2: Does discordance carry a return penalty?} Transaction-level OLS on 208{,}029 priced Buys (Equation~\ref{eq:rq2-ols}), with cluster-robust SE on \texttt{customerID} and two-sided Wald at $\alpha = 0.05$. Each customer contributes about 8 Buys on average, so within-customer residuals share unobserved shocks and motivate the clustered errors.
    \item \textbf{RQ3: Do baselines lift coherence on every band?} Per-customer PC@10 is averaged unweighted over (customer, split) within each (declared band, model) cell and per-band lift is the ratio of that mean to $\pi(b)$.
    \item \textbf{RQ4: Does PC-LightGCN improve coherence, and at what cost?} Aggregate metrics are averaged over the 69 splits, paired per-split deltas are computed against the baseline LightGCN, and the win rate is the share of splits on which the treatment beats the comparator.
\end{itemize}

The RQ2 OLS specification is
\begin{equation}
\label{eq:rq2-ols}
\begin{split}
r ={}& \beta_0 + \beta_1 \cdot \text{is\_coherent} \\
    &{}+ \beta_2 \cdot \text{asset\_volatility} \\
    &{}+ \mathbf{C}(\text{type})\boldsymbol{\gamma} \\
    &{}+ \mathbf{C}(\text{year})\boldsymbol{\delta} \\
    &{}+ \varepsilon.
\end{split}
\end{equation}
where $r = (p_{\text{end}} - p_{\text{start}})/p_{\text{start}}$ is the 6-month realised return as a decimal fraction (e.g., $r = 0.03$ means a 3 percentage-point gain), \texttt{is\_coherent} is binary ($d \leq 1$), and prices are matched via a backward-nearest join: each transaction's start price is the most recent close price on or before the trade date, and the end price is the most recent close price on or before the date six months later. Transactions without a price match within 7 days at either endpoint are dropped.

\subsection{Reproducibility}
\label{sec:setup:reproducibility}

All code, configuration grids, and evaluation artefacts are available in the project repository at \url{https://github.com/Samthesimpsons/SMU-Capstone}.
